\subsection{Recovering the full cutting group from a crossing arrow}

\begin{result}{E.29}{Proved}{The diagonal stabilizer has the full profinite topology}
Fix a stable labelled surface, its mapping class group $\Gamma$, and a
multicurve $\sigma$.  Put
\[
 C=\Stab_\Gamma(\sigma),\qquad
 \Delta=\langle\tau_c:c\in\sigma\rangle,\qquad H=C/\Delta.
\]
In the completed transporter exit category of A.4 and E.1, the
automorphism groups of the open and $\sigma$ objects are
$\widehat\Gamma$ and $\widehat H$, and
\[
 \Hom(\sigma,\varnothing)=\widehat\Gamma/\widehat\Delta.
                                                               \tag{E.79}
\]
The automorphism group of the distinguished crossing arrow, with
\emph{both} endpoints allowed to vary, is canonically
\[
 \widehat C\ \lhook\joinrel\longrightarrow\
       \widehat\Gamma\times\widehat H,
 \qquad c\longmapsto(c,\bar c).                                \tag{E.80}
\]
This is an embedding of the full profinite completion of $C$.
For the multicurve obtained by attaching a labelled three-holed
sphere to each boundary of a stable bordered surface, (E.80)
recovers the full completion of its boundary-fixed mapping class
group, including every boundary twist.
\end{result}
\begin{proof}
First $\Delta$ has its full profinite topology induced from
$\Gamma$.  The injection in \cite[Theorem 4.5]{ET_Boggi2014}
gives this already for the procongruence topology.  More precisely,
its proof on pages 17--19 uses the Looijenga levels to give, for
every $m\geq2$, a finite-index subgroup of $\Gamma$ whose
intersection with $\Delta$ is contained in $m\Delta$.
The first injection in that theorem follows from Theorem 2.1(ii)
and the level calculation cited there; it uses none of the
Section 5 centralizer claims audited in D.26.  Since
$\Delta\cong\Z^{|\sigma|}$, these subgroups give the asserted
full topology and an embedding $\widehat\Delta\to\widehat\Gamma$.

We prove the full topology assertion for the diagonal map
$C\to\Gamma\times H$ directly.  Let $K$ be any finite-index
subgroup of $C$.  Choose a finite-index normal subgroup
$L\triangleleft\Gamma$ with
$L\cap\Delta\subset K\cap\Delta$.  The subgroup
\[
 M=((K\cap L)\Delta)/\Delta\leq H
\]
has finite index.  The inverse image of the open product
$L\times M$ in $C$ is
\[
 L\cap(K\cap L)\Delta
       =(K\cap L)(L\cap\Delta)\subset K.
\]
Replacing $M$ by its core gives a finite quotient whenever one
is desired.  Every finite-index subgroup of $C$ therefore
contains an open subgroup induced from $\Gamma\times H$.
This proves (E.80) as a profinite embedding.  Completion of the
quotient sequence, together with the full topology of $\Delta$,
gives an exact sequence
\[
 1\longrightarrow\widehat\Delta\longrightarrow\widehat C
       \longrightarrow\widehat H\longrightarrow1.             \tag{E.81}
\]
For clarity, right exactness here follows by identifying all
finite quotients of $C$ that kill $\Delta$ with all finite
quotients of $H$; their common kernel is the closure of
$\Delta$ in $\widehat C$.

We next verify that these groups and the crossing set occur
\emph{simultaneously} in the strict categorical completion.
Use the finite skeletal set of multicurve types from A.4.
For a finite-index normal subgroup $N\triangleleft\Gamma$, let
$q:\Gamma\to Q=\Gamma/N$ and define a finite category with
the same objects by
\[
 \Hom_N(i,j)=q\bigl(\operatorname{Trans}(\sigma_i,\sigma_j)\bigr)
                         /q(\Delta_i).
                                                               \tag{E.82}
\]
These are subsets of the right-coset spaces $Q/q(\Delta_i)$.
For $g:i\to j$ and $h:j\to k$, composition is the class
of $hg$.  The transporter relation gives
$g^{-1}\Delta_jg\subset\Delta_i$.  Replacing $h$ by
$h\delta_j$ therefore changes $hg$ by the right factor
$g^{-1}\delta_jg\in\Delta_i$; replacing $g$ by
$g\delta_i$ has the same allowed effect.  The calculation
survives application of $q$, proving well-defined composition
and associativity in (E.82).

Every noninvertible arrow strictly decreases the number of
curves.  Consequently no composite involving a noninvertible
arrow is an endomorphism in this skeleton.  Each finite
endomorphism quotient in (E.82) can thus be refined by its
diagonal image with any prescribed finite quotient of the
original endomorphism group.  Keep the finite noninvertible
arrow sets and let the refined groups act through their old
quotients.  This gives a finite categorical quotient realizing
the prescribed endpoint quotients simultaneously with (E.82).
In particular all finite quotients of $\Gamma$ and $H$ occur.

Any finite categorical quotient restricts on the distinguished
crossing orbit to a finite transitive $\Gamma$-set.  Its action
kernel contains a finite-index normal subgroup $N$ of $\Gamma$,
so the quotient of $\Gamma/\Delta$ factors through
$\Gamma/(N\Delta)$.  Conversely (E.82) realizes each of these
sets.  Taking the inverse limit proves (E.79), with the full
endpoint actions just established.

At the discrete level those actions are
\[
 (g,\bar c)\cdot x\Delta=gxc^{-1}\Delta,
                  \qquad c\in C.
\]
They extend continuously.  Equation (E.81) identifies their
profinite homogeneous space with
\[
 (\widehat\Gamma\times\widehat H)/\widehat C
                 \ \cong\ \widehat\Gamma/\widehat\Delta:
                                                               \tag{E.83}
\]
lift the second coordinate through $\widehat C\to\widehat H$
to bring a coset to $\widehat\Gamma\times1$; two such
representatives differ precisely by $\widehat\Delta$.
This proves the asserted arrow stabilizer.  It does not require
the projection $\widehat C\to\widehat\Gamma$ to be injective.

Finally let the bordered core have type $(g,r)$ and attach one
three-holed sphere to each boundary, giving each sphere two
distinct external labels.  These labels make every graph
automorphism trivial and distinguish every branch of every
node.  The exact cutting-and-capping presentation is
\[
 C=\left(\Gamma_{g,r}^{\partial}\times(\Z^3)^r\right)
    \big/\left\langle
       t_i u_i^{-1},\ v_i,\ w_i\ (1\leq i\leq r)
    \right\rangle.                                           \tag{E.84}
\]
Here $t_i$ is the core boundary twist, $u_i$ the attached
inner boundary twist with compatible orientation, and
$v_i,w_i$ the two twists killed by external puncture capping.
The explicit map sends a core mapping class to its extension
and sends $u_i$ back to $t_i$, and $v_i,w_i$ to $1$.
These maps are inverse on the displayed generators and
relations.  Thus $C\cong\Gamma_{g,r}^{\partial}$, and (E.80)
recovers its full completion.  For $r=0$ the same statement
is the open-object automorphism group itself.
\end{proof}

The construction supplies the full-topology extraction map
required in E.23.  The separate question of recovering a whole
exit-system automorphism from its actions on these arrow
groups still includes its compatibility with all crossing
maps; that question is retained with its stated open status.
